\begin{textsample}{NEG Dim 6 – human – Score -3.00 – mg/mg\_ef\_ciencias\_5\_p2.txt}  \label{ex:f6_neg_004}
Nesse sentido, não basta que os conhecimentos científicos sejam apresentados aos estudantes. É preciso oferecer oportunidades para que eles, de fato, envolvam -se em processos de aprendizagem nos quais possam vivenciar momentos de investigação que lhes possibilitem exercitar e ampliar sua curiosidade, aperfeiçoar sua capacidade de observação, de raciocínio lógico e de criação ; desenvolver posturas mais colaborativas e sistematizar suas \textbf{primeiras} explicações sobre o mundo natural e tecnológico, e sobre seu corpo, sua saúde e seu bem-estar, tendo como referência os conhecimentos, as linguagens e os procedimentos próprios das Ciências da Natureza.
No que diz respeito, em especial, aos dois \textbf{primeiros} anos do Ensino Fundamental, em que se investe prioritariamente no processo de alfabetização das crianças, o ambiente pedagógico deve ser adequado para se desenvolver as habilidades de Ciências visando à progressão do \textbf{letramento} científico.

% matched lemmas: letramento, primeiro
\end{textsample}
